\subsection{Galois descent}

In this No.~we denote by N a field, \(\Gamma\) a group of automorphisms of N, \(\varepsilon\) the neutral element of \(\Gamma\) and K the field of invariants of \(\Gamma\) .

Let \(\mathbf{V}\) be a vector space over \(\mathbf{N}\) . We recall (II, p.~317) that a K-structure on \(\mathbf{V}\) is a vector sub-K-space \(\mathbf{V}_0\) of \(\mathbf{V}\) such that the K-linear mapping \(\varphi: \mathbf{N} \otimes_{\mathbf{K}} \mathbf{V}_0 \to \mathbf{V}\) which maps \(\lambda \otimes x\) to \(\lambda x\) is bijective. Let \(\mathbf{V}_0\) be such a K-structure; for each \(\sigma \in \Gamma\) we put \(u_\sigma = \varphi \circ (\sigma \otimes \mathrm{Id}_{\mathbf{V}_0}) \circ \varphi^{-1}\) ; then we have \(u_\sigma \left( \sum_{i \in I} \lambda_i e_i \right) = \sum_{i \in I} \sigma(\lambda_i) e_i\)

for every family of elements \(\lambda_{i}\) of \(\mathbf{N}\) and \(e_i\) of \(\mathbf{V}_0\) whence we obtain the relations

\[
u _ {\sigma} (x + y) = u _ {\sigma} (x) + u _ {\sigma} (y)\tag{1}
\]

\[
u _ {\sigma} (\lambda x) = \sigma (\lambda) u _ {\sigma} (x)\tag{2}
\]

\[
u _ {\sigma} \circ u _ {\tau} = u _ {\sigma \tau}\tag{3}
\]

\[
u _ {\varepsilon} = \operatorname{Id} _ {\mathrm{V}}\tag{4}
\]

for \(\sigma, \tau\) in \(\Gamma\) , \(x, y\) in \(\mathbf{V}\) and \(\lambda\) in \(\mathbf{N}\) .

PROPOSITION 6. --- a) Let V be a vector space over N with a K-structure. For a vector \(x \in V\) to be rational over K it is necessary and sufficient that \(u_{\sigma}(x) = x\) for all \(\sigma \in \Gamma\) . For a vector sub-N-space W of V to be rational over K it is necessary and sufficient that \(u_{\sigma}(W) \subset W\) for all \(\sigma \in \Gamma\) .

\begin{enumerate}
\def\labelenumi{\alph{enumi})}
\setcounter{enumi}{1}
\tightlist
\item
  Let \(V_{1}\) and \(V_{2}\) be two vector spaces over N, each with a K-structure. For a linear mapping of \(V_{1}\) into \(V_{2}\) to be rational over K it is necessary and sufficient that \(f(u_{\sigma}(x)) = u_{\sigma}(f(x))\) for all \(\sigma \in \Gamma\) and all \(x \in V_{1}\) .
\end{enumerate}

It is clear that K is the set of \(x \in N\) such that \(\sigma(xy) = x\sigma(y)\) for all \(\sigma \in \Gamma\) and all \(y \in N\) . The proposition therefore follows from Th. 1 (II, p.~324).

COROLLARY. --- Let \(V_{0}\) be a vector space over K and let W be a vector sub-N-space of \(N \otimes_{K} V_{0}\) . Suppose that W is stable under the mappings \(\sigma \otimes Id_{V_{0}}\) for all \(\sigma \in \Gamma\) . Let \(W_{0}\) be the set of \(x \in V_{0}\) such that \(1 \otimes x \in W\) ; then \(W_{0}\) is the unique vector sub-K-space of \(V_{0}\) such that \(W = N \otimes_{K} W_{0}\) .

It suffices to remark that the set of elements of the form \(1 \otimes x\) ( \(x \in \mathbf{V}_0\) ) is a K-structure on \(\mathbf{N} \otimes_{\mathbf{K}} \mathbf{V}_0\) for which we have \(u_\sigma = \sigma \otimes \mathrm{Id}_{\mathbf{V}_0}\) for \(\sigma \in \Gamma\) .

PROPOSITION 7. --- Let V be a vector space over N, \((u_{\sigma})_{\sigma \in \Gamma}\) a family of mappings of V into itself satisfying (1) to (4) and \(V_{0}\) the set of \(x \in V\) such that \(u_{\sigma}(x) = x\) for all \(\sigma \in \Gamma\) .

\begin{enumerate}
\def\labelenumi{\alph{enumi})}
\item
  \(\mathbf{V}_0\) is a vector sub-K-space of \(\mathbf{V}\) and the K-linear mapping \(\varphi\) of \(\mathbf{N} \otimes_{\mathbf{K}} \mathbf{V}_0\) into \(\mathbf{V}\) which maps \(\lambda \otimes x\) to \(\lambda x\) is injective.
\item
  If \(\Gamma\) is finite, then \(\varphi\) is bijective and \(\mathbf{V}_0\) is a K-structure on \(\mathbf{V}\) .
\end{enumerate}

It is clear that \(\mathbf{V}_0\) is a vector sub-K-space of \(\mathbf{V}\) .

The formula \(u_{\sigma} \circ \varphi = \varphi \circ (\sigma \otimes \operatorname{Id}_{\mathbf{V}_{0}})\) shows that the kernel W of \(\varphi\) is stable under the mappings \(\sigma \otimes \operatorname{Id}_{\mathbf{V}_{0}}\) ; by the Cor. to Prop. 6 there exists therefore a subspace \(W_{0}\) of \(V_{0}\) such that \(W = N \otimes_{K} W_{0}\) . If x belongs to \(W_{0}\) we then have \(x = \varphi(1 \otimes x) = 0\) , hence \(W_{0} = 0\) and so W = 0. This proves a).

Suppose now that \(\Gamma\) is finite; we have to show that \(\varphi\) is surjective, or equivalently that \(V_0\) generates the vector N-space V. Thus let \(f\) be an N-linear form on V whose restriction to \(V_0\) is zero. Let \(x \in V\) ; for every \(\lambda \in N\) the element \(y_\lambda = \sum_{\sigma \in \Gamma} u_\sigma(\lambda x)\) of V clearly belongs to \(V_0\) , whence \(f(y_\lambda) = 0\) , that is,

\(\sum_{\sigma \in \Gamma}f(u_{\sigma}(x))\sigma (\lambda) = 0\) . By Dedekind's theorem (V, p.~27, Cor. 2) we thus have

\(f(u_{\sigma}(x)) = 0\) for each \(\sigma \in \Gamma\) ; in particular, taking \(\sigma = \varepsilon\) we find \(f(x) = 0\) , which means that \(f = 0\) . This proves \(b\) ).

Let \(\mathbf{M}\) be a vector space over \(\mathbf{N}\) ; for each \(\sigma \in \Gamma\) let \(\mathbf{M}^{\sigma}\) be the vector space over \(\mathbf{N}\) with the same underlying additive group as \(\mathbf{M}\) , with the external law \((\lambda, x) \mapsto \sigma(\lambda)x\) . Write \(\mathbf{V} = \prod_{\sigma \in \Gamma} \mathbf{M}^{\sigma}\) ; the underlying additive group of \(\mathbf{V}\) is that of

all mappings of \(\Gamma\) into \(\mathbf{M}\) , with the external law defined by

\[
(\lambda . h) (\sigma) = \sigma (\lambda) h (\sigma) \quad (\lambda \in \mathrm{N}, h \in \mathrm{V}, \sigma \in \Gamma).\tag{5}
\]

(The product \(\sigma(\lambda) h(\sigma)\) is calculated in the vector space M.) Further, we define on \(\mathbf{N} \otimes_{\mathbf{K}} \mathbf{M}\) a vector space structure over \(\mathbf{N}\) by the formula

\[
\lambda \left(\sum_ {i} \mu_ {i} \otimes x _ {i}\right) = \sum_ {i} \lambda \mu_ {i} \otimes x _ {i}.
\]

Finally we denote by \(\psi\) the K-linear mapping of \(\mathbf{N} \otimes_{\mathbf{K}} \mathbf{M}\) into \(\mathbf{V}\) characterized by the relation

\[
\psi (\lambda \otimes x) (\sigma) = \sigma (\lambda). x\tag{6}
\]

for \(\lambda \in \mathbf{N}\) , \(x \in \mathbf{M}\) and \(\sigma \in \Gamma\) . It is clear that \(\psi\) is N-linear.

PROPOSITION 8. --- The N-linear mapping \(\psi\) of \(\mathbf{N} \otimes_{\mathbf{K}} \mathbf{M}\) into \(\mathbf{V} = \prod_{\sigma \in \Gamma} \mathbf{M}^{\sigma}\) is

injective, and it is bijective if \(\Gamma\) is finite.

For every \(\sigma\in\Gamma\) we define a mapping \(u_{\sigma}\) of V into V by

\[
\left(u _ {\sigma} h\right) (\tau) = h (\tau \sigma)\tag{7}
\]

for \(h \in \mathbf{V}\) and \(\tau \in \Gamma\) . The verification of (1)-(4) is immediate. Denote by \(\mathbf{V}_0\) the set of \(h \in \mathbf{V}\) such that \(u_{\sigma}(h) = h\) for all \(\sigma \in \Gamma\) . For each \(x \in \mathbf{M}\) let \(\theta(x)\) be the constant mapping of \(\Gamma\) into \(\mathbf{M}\) with value \(x\) ; then \(\theta\) is a K-isomorphism of \(\mathbf{M}\) onto \(\mathbf{V}_0\) . If we define the homomorphism \(\varphi: \mathbf{N} \otimes_{\mathbf{K}} \mathbf{V}_0 \to \mathbf{V}\) as above, we have \(\psi = \varphi \circ (\mathrm{Id}_{\mathbf{N}} \otimes \theta)\) and now Prop. 8 follows from Prop. 7.

COROLLARY. --- Let \(\psi\) be the K-linear mapping of N \(\otimes_{K}\) N into the product vector space \(N^{\Gamma}\) such that \(\psi(x \otimes y) (\sigma) = \sigma(x)y\) for x, y in N and \(\sigma \in \Gamma\) . Then \(\psi\) is injective and it is bijective when \(\Gamma\) is finite.

This is the particular case \(\mathbf{M} = \mathbf{N}\) of Prop. 8.

Remarks. --- 1) Let F be an extension of K and N a subextension of F, and let \(\Gamma\) be a finite group of automorphisms of N, of which K is the field of invariants. Prop. 8 implies the existence of an isomorphism of K-algebras \(\theta: \mathbf{N} \otimes_{\mathbf{K}} \mathbf{F} \to \mathbf{F}^{\Gamma}\) characterized by \(\theta(x \otimes y)\) ( \(\sigma\) ) = \(\sigma(x)y\) for \(x \in \mathbf{N}\) , \(y \in \mathbf{F}\) and \(\sigma \in \Gamma\) .

\begin{enumerate}
\def\labelenumi{\arabic{enumi})}
\setcounter{enumi}{1}
\tightlist
\item
  The notation K, N and \(\Gamma\) has the same meaning as before. For each integer \(n \geqslant 1\) , let \(A_{n}\) be the tensor product of n K-algebras identical with N; let \(B_{n}\) be the set of mappings of \(\Gamma^{n-1}\) into N. By induction on n we can deduce from the Cor. of Prop. 8 the existence of an isomorphism \(\varphi_{n}: A_{n} \to B_{n}\) mapping \(x_{1} \otimes \ldots \otimes x_{n}\) to the function \((\sigma_{1}, \ldots, \sigma_{n-1}) \mapsto \sigma_{1}(x_{1}) \ldots \sigma_{n-1}(x_{n-1}) x_{n}\) .
\end{enumerate}

